\section {Formal verification overview}
Formal Methods refers to mathematically rigorous techniques and
tools for the specification, design and verification of software and
hardware systems.\\ 
Expanding this definition Formal methods 
can be used for formal description of the system (specification) 
on some level of abstraction, program synthesis (design) according to specification
and verification of a system to prove that model of system under consideration 
satisfies some properties.

%Let us get started by citing \cite{BeckKKF13}!
%So what did \citeauthor{BeckKKF13} do in \citeyear{BeckKKF13}?
%Maybe it is answered in \cref{introduction} on page~\pageref{introduction}?

% You can look it up on \url{https://en.wikipedia.org/wiki/Pythagorean_theorem}.
Formal verification techniques can be thought of as comprising three parts \cite{HuthR1}:
\begin{itemize}
    \item a framework for modelling systems, typically a description language;
    \item a specification language for describing the properties to be verified;
    \item a verification method to establish whether the description of a system satisfies the specification.
\end{itemize}
These methods and techniques of verification include:
\begin{enumerate}

    \item SAT solving - operate on propositional logic 

    \item SMT solving (and tools like Z3) extend this to first-order logic,
        integrating theories such as arithmetic or arrays.

    \item Relational model finding (e.g., Alloy) 

    \item Model checking (e.g., NuSMV) exhaustively explore state-space against CTL/LTL specifications.

    \item Interactive theorem provers (e.g., Coq/Rocq, Isabelle/HOL) 
\end{enumerate}


Approaches to verification can be classified according to the following criteria \cite{HuthR1} pages 172-173:\\
\textbf{Proof-based vs. model-based}\\
In a proof-based approach, the system description 
is a set of formulas $\Gamma$ (in a suitable logic) and the specification is another formula $\phi$. 
The verification method consists of trying to find a proof that $\Gamma \vdash \phi$. 
In a model-based approach, the system is represented by a model M for an appropriate logic.
The specification is again represented by a formula $\phi$ and the verification method consists 
of computing whether a model M satisfies $\phi$ ($M \models \phi$). This computation is usually automatic
for finite models. Logical proof systems are often sound and complete, meaning that $\Gamma \vdash \phi$ (provability) holds if, and only
if, $\Gamma\models\phi$(semantic entailment) holds, where the latter is defined 
as follows: 
for all models M, if for all $\psi \in \Gamma$ we have M $\models\psi$, then M $\models\phi$.
Thus, we see that the model-based approach is potentially simpler than
the proof-based approach, for it is based on a single model M rather
than a possibly infinite class of them.

\textbf{Degree of automation. }\\ 
Approaches differ on how automatic the method is; 
the extremes are fully automatic and fully manual. Many
of the computer-assisted techniques are somewhere in the middle.

\textbf{Full- vs. property-verification.}\\ 
The specification may describe a single property of the system, or it may describe its full behavior. 
The latter is typically expensive to verify.
Intended domain of application, which may be hardware or software;
sequential or concurrent; reactive or terminating; etc. A reactive system
is one which reacts to its environment and is not meant to terminate
(e.g., operating systems, embedded systems and computer hardware).

\textbf{Pre- vs. post-development. }\\ 
Verification is of greater advantage if introduced early in the course of system development, because errors
caught earlier in the production cycle are less costly to rectify. (It is
alleged that Intel lost millions of dollars by releasing their Pentium chip with the FDIV error.


\subsection {SAT solvers}

Propositional logic - consider  propositions and relations between them and 
constructions of argument based on them.

Propositions - facts or statements about the world  that can be true or false
. Also defined as atomic propositions or Atoms.
Some propositions can be described in natural language as atomic or indecomposable 
declarative sentences) \cite{HuthR1}.


Atomic propositions can be connected to each other with logical connectives and compose propositional formulas.
The most common connective relations are:
\begin{itemize}
    \item Negation $ \neg $
    \item Conjunction $ \land $ 
    \item Disunction $ \lor $ 
    \item Implication $ \rightarrow $
\end{itemize}

The Backus Naur form (BNF) for propositional logic well-formed formula:
\[ \phi ::= p\ |\ (\neg \phi)\ |\ (\phi \land \phi)\ |\ (\phi\lor\phi)\ |\ (\phi\rightarrow\phi)    \]
where p - is any atomic proposition.

Based on these we can construct a system of reasoning: from some set of formulas we can infer (prove) another formula that can be true or false:
\[ \phi_1, \phi_2, \phi_3, ... \phi_n \vdash \psi \]
Such formulas called sequents, where on the left side there is a set of formulas called premises and on the right side - a formula called conclusion


Another relation between logical formulas  is semantic entailment:
\[ \phi_1, \phi_2, \phi_3, ... \phi_n \models \psi \]
it means that if propositional logic formulas $\phi_1 .. \phi_n$ evaluate to True, then formula $\psi$ also evaluates to True.

The formula $\phi$ is valid if for all valuations of atomic propositions in formula it valuates to True:
\[ \models \phi \]


SAT solver is a program which aims to solve the Boolean satisfiability problem 
(SAT), where the Boolean satisfiability problem sometimes called 
propositional satisfiability asks whether there 
exists an interpretation that satisfies a given Boolean formula. 

Given formula  $\phi$ is \textbf{satisfiable} if it has a valuation in which it valuates to True.:

In other words, it asks whether the formula's variables can be 
consistently replaced by the values TRUE or FALSE to make the formula 
evaluate to TRUE. If this is the case, the formula is called satisfiable, 
else unsatisfiable.\cite{HuthR1}.

SAT problem is NP-complete, as a result, only algorithms with exponential 
worst-case complexity are known. In spite of this, efficient and 
scalable algorithms for SAT were developed during the 2000s, which have 
contributed to dramatic advances in the ability to automatically 
solve problem instances involving tens of thousands of variables and 
millions of constraints.%[https:en.wikipedia.org/wiki/SAT_solver -1.]

SAT solvers usually begin by converting a problem to conjunctive normal form (CNF).
\textbf{Conjunctive normal form (CNF)} - as a conjunction of clauses, where each clause is a disjunction of literals.
Literal is an atomic proposition or a negated atomic proposition.
Conjunctive normal form allows easy check of validity of a formula which otherwise take times exponential to number of atoms.




% \subsection {SMT}
% Propositional logic is limited to express relations between logical components other than equivalent to natural language: \emph{not, and, or, if..then} and 
% declarative sentences that can only be True or False. 
%
% In order to be able to operate with finer logical structures we need to enrich the language.
% The necessity to solve this problem lead to design of Predicate logic or first-order logic.\\
% In First-order logic we introduce:\\
% \begin{itemize}
%     
%     \item Variables that represent individual objects.
%     \item Functions $\mathcal{F}$ that refer to objects.
%     \item Predicates $\mathcal{P}$ - symbols that semantically represents some property or relation.
%     \item Quantifiers: universal $\forall$ and existential $\exists$.
% \end{itemize}
%
%
% The Backus Naur form (BNF) for predicate logic formulas:
%
% \[ \phi ::= P(t_1,t_2,...t_n)\ |\ (\neg \phi)\ |\ (\phi \lor \phi)\  |\ (\phi \land \phi)\   |\ (\phi \rightarrow \phi)\ |\ (\forall x \phi)\ |\ (\exists x \phi)\ \]
%
% where,\\ 
% t - called Terms and can be variable, constants or functions on terms  \[ t ::= x\ |\ c\ |\ f(t,...t) \]
%
% Atomic formulas are predicates applied to terms: $P(t_1,t_2,...t_n) $
%
% Variables can be free (supplied from outside the formula) or bounded by some quantifier.
%
% Unlike propositional logic, in predicate logical we cannot just evaluate formulas based on given valuation by assigning True or False to the variables.  
% We require a model of all function and predicate symbols involved.\\ 
% A model M for a first‑order logic contains: 
% \begin{itemize}
%     \item Domain A is a non‑empty set of individuals, mapping each variable to a value in A
%     \item Set of functions where n-ary function is interpreted as a concrete function $f^m : A^n \rightarrow A$.
%     \item Set of predicates n-ary predicate is interpreted as a relation $P^m \subseteq A^n $ 
% \end{itemize}
%
% Truth in a model is defined recursively:
%
% \[M \models P(t_1,...,t_n) iff (t_1,...,t_n) \in P^m \]
%
% % TODO!!! formula and citations
% \[M\models\ \text{ forall }\ x\phi\ \text{iff for every a} \ \in A, M \models [x \rightarrow a ] \phi \]
% \[M\models\ \text{exists} \ x\phi \ \text{iff for some a} \  \in A, M \models [x \rightarrow a ] \phi \]
%
% A sentence has no free variables, so its truth value depends solely on M.  
%
% We write $\Gamma\models\psi$ (semantic entailment) when every model that makes all formulas in $\Gamma$ true also makes $\psi $ true. 
%
% Soundness and completeness ensure $\Gamma\vdash\psi$ (provability) exactly when $\Gamma\models\psi$,
% but the two perspectives complement each other: proofs are constructive, while counter‑models establish non‑entailment.
%


\section{Model checking}

Model checking is based on temporal logic. 
The model in temporal logic described as a system with few states that can change with time.
The formula is not statically true or false in a model, as it is in propositional and predicate logic,
it can be true in some states and false in others, the static notion is replaced 
by a dynamic one \cite{HuthR1}.

In model checking, the models M are transition systems and the properties $\phi$ are formulas in temporal logic. 
To verify that a system satisfies a property, we must do three things:
\begin{enumerate}
    \item model the system using the description language of a model checker, arriving at a model M;
    \item code the property using the specification language of the model checker, resulting in a temporal logic formula $\phi$;
    \item Run the model checker with inputs M and $\phi$.
\end{enumerate}



Model checking is a core approach in Formal verification that involves an algorithmic search of the state space of a system model to check if a
given property (specification) holds, where:\\
Model- is an abstract representation of the system (often a state-transition graph or a
program) capturing its possible states and transitions. The model can be described in a modeling language
(like Promela for SPIN, or as a transition system in a tool like NuSMV). \\
Specification- is a formal property that the model should satisfy, commonly expressed in a temporal logic such as LTL or CTL. For example, an LTL
specification might state that “whenever request happens, eventually response occurs” or a CTL property
might require that “in all execution paths, a certain error state is never reached.” \\ 


\subsection {Classification of Model checking methods}
Model checking can be classified based on time concept into:
\textbf{linear-time temporal logic } and \textbf{branching time logic.}
\cite{HuthR1}
\\


\subsubsection{Linear-time temporal logic }
Linear-time temporal logic (LTL) considers time as a set of paths, where each path 
is a sequence of time instances \cite{HuthR1}.
LTL has the following syntax in Backus Naur form:
\[ \phi ::= T\ |\  \perp |\  p\ |\  (\neg\phi)\  |\  (\phi \land \phi)\  |\  (\phi \lor \phi) |\ (\phi \rightarrow \phi)\   \]
\[|\   (X \phi)\  |\  (F \phi)\  |\  (G \phi)\  |\  (\phi\ U\ \phi)\  |\  (\phi\ W\ \phi)\  |\  (\phi\ R\ \phi)\ \]
The connectives $X, F, G, U, W, R$ called temporal connectives, because they describe relations between states in different moments of time.
The meaning of these connectives described in Table \ref{tab:LTL1}\\

\begin{table}[h!]   
    \begin{center}
    \begin{tabular}{| m{0.6cm} | m{3.5cm} | m{9.5cm} |}
        \hline
        X & Next & $\phi$ must hold in the very first successor state on the chosen path.\\
        \hline
        F & Finally ($F \phi$) & $\phi$ must hold at some future state on the path  (i.e., “eventually $\phi$”).\\
        \hline
        G & Globally ($G \phi$) & $\phi$ must hold at every state on the entire path  (“always $\phi$”).\ \ \ \ \ \ \ \ \ \ \\
        \hline
        U & Until ($\phi U \psi$) & $\psi$ must be true at a future state, $\phi$ must stay true at every state before that one.\\
        \hline
        W & Weak until ($\phi W \psi$) & either ($\phi U \psi$) holds or $\phi$ stays true forever if $\psi$ never becomes true.\\
        \hline
        R & Release ($\phi R \psi$) & $\psi$ must hold up to and including the first state where 
        $\phi$ becomes true; if $\phi$ never becomes true, $\psi$ must hold forever (the dual of Until).\\
        \hline
    \end{tabular}
    \caption{LTL temporal connectives}
\end{center}
\end{table}
    \label{tab:LTL1}

In LTL, model can be represented as a transition system \cite{DGL1}.\\
Transition system $T$:
\[ T = (S,\stackrel{a}{\rightarrow}, L,s) \] 
where:\\
S - set of states, called the space states of T,\\
$\stackrel{a}{\rightarrow}$ - transitions (binary relation of states), associated with some action $a \in A$\\
L - labelling function that reflects states to propositional atoms\\
s - initial state of T, can be omitted if we want to describe only global structure.\\
Sometimes if we consider only one action label we can omit labeling transition: $\rightarrow$.\\
A path $\pi$ in a transition system $T$ is a finite or infinite sequence of states and (optionally) actions
which transform each state into its successor: $s_0 \stackrel{a_0}{\rightarrow} s_1 \stackrel{a_1}{\rightarrow} ... $

Trace or computation in a transition system is a finite or infinite sequence of state descriptions along the path:
$L(s_0) \stackrel{a_0}{\rightarrow} L(s_1) \stackrel{a_1}{\rightarrow} ... $

\subsubsection{Branching temporal logic (CTL)}
Branching temporal logic represents time as a tree rooted in present and 
branching into the future, the other name is Computation Tree logic (CTL).
LTL implicitly quantifies universally over paths \cite{DGL1}.  Therefore, properties which 
assert the existence of a path cannot be expressed in LTL. 
We can solve this problem by considering the negation of the property
in question. For example to check whether there exists a path from s satisfying the LTL formula $\phi$, we check whether all paths
satisfy $\neg\phi$; a positive answer to this is a negative answer to our original question, and vice versa.
We used this approach when analysing the ferryman puzzle in the previous section. 
However, properties which mix universal and existential path quantifies cannot in general be model
checked using LTL approach, because the complement formula still has a mix.
Branching-time logics solve this problem by allowing us to quantify explicitly over paths.

CTL has the following syntax in Backus Naur form:
\[ \phi ::= T\ |\  \perp |\  p\ |\  (\neg\phi)\  |\  (\phi \land \phi)\  |\  
(\phi \lor \phi) |\ (\phi \rightarrow \phi)\   \]

\[|\   (AX \phi)\  |\  (EX \phi)\  |\  (AF \phi)\  |\  (EF \phi)\  |\  
(AG \phi)\  |\  (EG \phi)\  |\  A[\phi\ U\ \phi]\  |\  E[\phi\ U\ \phi]\  \]

New introduced quantifiers A and E mean:\\
$A \phi $ - formula $\phi$ for all paths \\
$E \phi $ - exists a path that formula $\phi$\\
Connectives X,F,G cannot occure without being preceded by A or E.\\
Weak until W and release R are usually not included in CTL, but can be derived.


\subsubsection{Temporal logic (CTL*)}
\begin{figure}[!h]
    \centering
    \includegraphics[width=12cm]{CTL} 
    \caption{LTL, CTL and CTL* logic relation \cite{HuthR1}}
    \label{fig:CTL}
\end{figure}


\subsubsection{Statistical Model Checking (SMC)}
Some systems have huge or continuous state spaces (e.g. complicated analog components, 
or software with random timing), or they exhibit probabilistic behavior 
(randomized protocols, hardware failures) where one is interested in quantitative properties (like
"the request will be served within 1s with probability at least 0.99"). 
Statistical Model Checking (SMC) is an approximate verification technique that uses simulation 
and statistical methods instead of exhaustive exploration \cite{Younes2002AcceptanceSampling}. 

SMC treats the model checking problem as a statistical inference problem \cite{Younes2006StatisticalPMC}.
It executes a number of simulations (random or guided) of the system. 
Checks whether each simulation trace satisfies the property $\phi$ 
(which often is a bounded temporal logic property like “within T time units,
something happens”). 
Based on the outcomes, it uses statistics (like hypothesis testing or estimation with
confidence intervals) to infer with high confidence whether the property holds with required probability.

SMC was pioneered by Younes in the 2000s \cite{Younes2002AcceptanceSampling}. 
There are two main modes: hypothesis testing (e.g. "check if the probability 
$P($\textit{property}$) > p$ with confidence 95\%") and estimation (estimate the
probability itself within some error bound). 

They adopt the continuous stochastic logic (CSL) as our formalism for expressing 
probabilistic real-time properties of discrete event systems. CSL has previously 
been proposed as a formalism for expressing temporal and probabilistic
properties of continuous-time Markov chains (CTMCs) by \citeauthor{Aziz1996VerifyingCTMC} 
\cite{Aziz1996VerifyingCTMC}.

A CSL formula is either a state formula or a path formula.

The Backus Naur form (BNF) for CSL logic well-formed formula:

State formulas:
\[
    \phi ::= \text{true}\ |\ a\ |\ \neg \phi\ |\ \phi \land \phi\ |\ P_{\leq \theta} (\rho)(\phi)
\]

Path formulas:
\[
    \rho::= X \phi\ |\ \phi U_{\leq t} \phi
\]
Where $\phi$ is a state formula and  $\rho$ is a path formula.

Tools like PRISM and its SMC extension, described in \cite{Kwiatkowska2011PRISM}, 
STORM \cite{Dehnert2017Storm}, or PLASMA Lab \cite{Legay2016PlasmaLab} implement 
statistical model checking for various kinds of stochastic systems.

The obvious advantage is that SMC sidesteps the state explosion by not exploring the full
state space – instead, it samples paths. 
This makes it scalable to systems with enormous or even infinite state
spaces, as long as we can simulate the system. It is also applicable to systems 
with continuous dynamics or very large numbers of states where traditional model 
checking is infeasible . Another benefit is ease of implementation: if there 
is a simulator for your system, adding a statistical checker on top is easier 
than building a full model checker. SMC naturally supports probabilistic properties 
– e.g., verifying properties of a Markov Chain or Markov Decision Process by 
sampling behaviors and using results like Chernoff or Hoeffding bounds to determine 
if the property likely holds. Because it uses actual simulations, SMC can
handle very complex nonlinear or black-box system aspects (treating them as 
black boxes that just produce simulation traces).

The trade-off is that SMC provides a confidence result, not a guarantee. It is probabilistic
verification: it might say “with 99\% confidence the property holds with probability $\leq 0.97$”. There’s a small
chance of error in the verification result (false positives/negatives bounded by the confidence level). Also,
SMC may require a large number of simulations to reach high confidence, especially for rare events (if the
property failing is a rare event, one might need many samples to observe it even once). Rare property
verification is a known challenge – advanced importance sampling or rare event simulation techniques
might be needed in such cases. 

SMC can’t truly prove a property in the mathematical sense – it can only give statistical evidence.
If a property almost always holds but has a corner-case failure, SMC
might miss it if that scenario has very low probability and isn’t sampled. Thus, it sacrifices completeness.
Another limitation is that SMC typically handles quantitative properties (often expressed in probabilistic
temporal logics like PCTL or BLTL) – it’s not usually used for pure logical properties that are either true or
false; it’s more for properties like “the probability of an event is at least p”.



\subsubsection{Runtime Verification and Monitoring}

Runtime verification (RV) is a technique that involves monitoring a running 
system (or execution traces produced by it) against a formal specification \cite{Hinrichs2014ModelCheckRuntimeVerify}. 
Unlike the other methods discussed, which are largely offline
(analyzing a model mathematically), runtime verification is applied during execution. It can be seen as a
lightweight formal method, sitting somewhere between testing and model checking: you instrument the
system with monitors, and these monitors check whether the execution is satisfying the desired properties.
If a violation is detected, it can be logged, or in some cases, a recovery action can be taken.

Key components of runtime verification include:\\ 
- A set of monitors or runtime assertions derived from the formal specification 
(often temporal logic formulae or automata on traces).\\
- An instrumentation of the system to report events to the monitors.\\
- The monitors run in parallel with the system (or on logged traces)\\
and detect property satisfaction or violation on the fly.

For example, if you have a property “whenever interrupt is raised, service routine runs within 100ms,” a
runtime monitor can watch for interrupts and check timing until service routine is called, flagging an error if
the 100ms deadline is missed.

Runtime verification avoids the state explosion problem entirely, because it doesn’t explore all
behaviors — it only checks the actual behaviors that occur. Runtime verification 
avoids the complexity of traditional formal verification techniques (like model 
checking and theorem proving) by analyzing only one or a few execution traces 
and by working directly with the actual system. This means it scales well to 
very complex systems, since you’re just running the system (or a
high-fidelity simulation of it) and the overhead is only in checking the property 
on that run. It also has the advantage of checking the real system (no abstraction 
or modeling error, since the code that runs is what’s verified against the property
, at least for that run). It can catch errors in scenarios that were not anticipated
by designers as soon as they occur in the field or during testing, thus serving 
as a powerful debugging and validation aid. Runtime verification can be used 
continuously in deployed systems for critical properties
(e.g., a monitor that ensures no unauthorized file access operations occur, 
and triggers an alarm if they do, providing a layer of security monitoring).
Another plus is that runtime verification can sometimes detect issues that static model checking might miss
due to environment assumptions. At runtime, the real environment provides inputs and the monitor can
see violations under real conditions, including those hard-to-model aspects.


The fundamental limitation is that runtime verification is not exhaustive. It 
can only inspect the executions that actually happen (or a sample of them). 
If a particular bug does not manifest in the observed runs, RV cannot detect 
it. It provides no guarantee that “no bug exists” — only that “no bug was
observed in the runs we monitored”. In other words, it lacks completeness/coverage
, unlike model checking which (for a model) covers all behaviors by design. 
Another challenge is performance overhead:
monitors consume resources, and if properties are complex (say, specified in 
LTL, which might require monitoring a suffix of the trace), the runtime checking 
might slow the system or require buffering events.
However, in many cases monitors can be optimized or run on a separate core, 
and the overhead is manageable (reports of a few percent overhead for simple 
monitors are common).
Additionally, writing good monitors (formal properties) and instrumenting the 
system requires effort. It’s easier than full verification, but still requires 
formal specification skills. If one instrument too much or
monitors too many properties, the overhead can become significant, or the system
’s real-time behavior might be perturbed.


Runtime verification is often used in safety-critical systems as an additional 
runtime safety net. 
For instance, in aerospace or automotive software, one might include monitors 
for conditions like “no two actuators turn on simultaneously” 
or “if sensor reading exceeds threshold , system enters safe mode
within 1 second,” etc. NASA has explored RV for spacecraft software monitoring 
during missions \cite{Slagel2024FormalRuntimeAssurance}(since you cannot model check the entire flight code with all 
analog intricacies, but you can monitor key properties as
the mission progresses). Runtime Verification has also become popular in security 
(e.g., monitoring program execution for compliance with security policies). 
For example, a monitor can watch system calls to detect a sequence that matches 
a known attack pattern or a violation of a security automaton.
Another domain is testing: during testing, using runtime monitors can greatly 
increase the chance of noticing an issue. Instead of manually inspecting outputs
, formal monitors can rigorously check that each test execution satisfies the 
expected temporal properties. This is often called “monitoring oracles” in testing
– formal specifications serve as test oracles at runtime.
In summary, runtime verification does not replace model checking but complements 
it. Where model checking cannot be applied (due to scale or undecidability), 
or in deployed scenarios where you want continual assurance, RV provides a way 
to leverage formal specifications dynamically. It trades completeness for scalability 
and immediacy.

% \subsection {Comparison of Model checking methods}


% State explosion - the exponential growth of the state space with system parameters (like number of components or data bits).
% Combating state explosion is a central theme in model checking research


% \subsection {Application of Model checking}
% \subsubsection{Hardware model checking}
% \subsubsection{Software model checking}
% \subsubsection{Reactive systems model checking}
%
% Reactive system: a system that maintains an
% ongoing interaction with its environment, responding to inputs with outputs continually (examples include
% operating systems, web servers, control systems). Reactive systems are typically non-terminating and
% concurrency-intensive. 




